\section{Runtime Semantics}
\label{sec:runtimesemantics}

This appendix provides the formal runtime semantics for the AxiomQuant Five
Redux operator chain. The semantics define how the operators $\Phi_5, \Phi_7,
\Phi_8, \Phi_8^t, \Phi_9$ are implemented in the Julia runtime while preserving
the mathematical properties required by the canonical spine.

\subsection{Runtime Model}

Let the runtime state be:
\[
\mathcal{R}(x) = (g(x), R(x), H(x), \tau(x), s(x), c(x), h(x), t(x), I(x)).
\]

The runtime must satisfy:
\[
\mathcal{R}(x) \in C^\infty(\mathcal{M})
\]
and preserve altitude/domain semantics.

\subsection{Operator Contracts}

Each operator has a runtime contract:

\[
\Phi_k^{\text{rt}} : \text{State}_{k-1} \to \text{State}_k,
\]
where:
\[
\text{State}_0 = x,
\qquad
\text{State}_1 = s(x),
\qquad
\text{State}_2 = c(x),
\qquad
\text{State}_3 = h(x),
\qquad
\text{State}_4 = t(x),
\qquad
\text{State}_5 = I(x).
\]

Runtime operators must satisfy:
\[
\Phi_k^{\text{rt}} = \Phi_k + \epsilon_k,
\]
where $\epsilon_k$ is a bounded numerical error.

\subsection{Runtime Smoothness}

Each operator must satisfy:
\[
\Phi_k^{\text{rt}} \in C^\infty,
\]
and numerical approximations must preserve smoothness:
\[
\| \epsilon_k \| < \delta_k,
\qquad
\delta_k \ll 1.
\]

\subsection{Runtime Stability}

Let $C_k$ be the stability constant for $\Phi_k$.

Runtime stability requires:
\[
\| \Phi_k^{\text{rt}}(u) - \Phi_k^{\text{rt}}(v) \|
\le
C_k \| u - v \| + \delta_k.
\]

Thus the full pipeline satisfies:
\[
\| E^{\text{rt}}(x) - E^{\text{rt}}(x') \|
\le
C \| x - x' \| + \delta,
\]
where:
\[
C = C_9 C_8^t C_8 C_7 C_5,
\qquad
\delta = \sum_k \delta_k.
\]

\subsection{Runtime Reversibility}

Runtime reversibility requires:
\[
E^{\text{rt}}(x) = E^{\text{rt}}(x')
\Rightarrow
x = x' + \eta,
\]
where $\eta$ is bounded numerical error:
\[
\| \eta \| < \delta.
\]

\subsection{Runtime Monotonicity}

Runtime monotonicity requires:
\[
u \le v \Rightarrow \Phi_k^{\text{rt}}(u) \le \Phi_k^{\text{rt}}(v).
\]

Thus:
\[
x \le x' \Rightarrow E^{\text{rt}}(x) \le E^{\text{rt}}(x').
\]

\subsection{Runtime Altitude Preservation}

Runtime altitude preservation requires:
\[
A_k(u) \le A_k(v)
\Rightarrow
A_{k+1}(\Phi_k^{\text{rt}}(u))
\le
A_{k+1}(\Phi_k^{\text{rt}}(v)).
\]

This ensures the runtime pipeline respects the altitude ladder.

\subsection{Runtime Domain Preservation}

Runtime domain preservation requires:
\[
u \in D_k \Rightarrow \Phi_k^{\text{rt}}(u) \in D_{k+1}.
\]

Thus:
\[
x \in D_1 \Rightarrow E^{\text{rt}}(x) \in D_5.
\]

\subsection{Runtime Pipeline}

The runtime pipeline is:
\[
E^{\text{rt}} = 
\Phi_9^{\text{rt}}
\circ
\Phi_8^{t,\text{rt}}
\circ
\Phi_8^{\text{rt}}
\circ
\Phi_7^{\text{rt}}
\circ
\Phi_5^{\text{rt}}.
\]

Runtime correctness requires:
\[
E^{\text{rt}} = E + \epsilon,
\qquad
\| \epsilon \| < \delta.
\]

\subsection{Operator Role}

The runtime semantics ensure that the executable Julia pipeline preserves:

\begin{itemize}
    \item smoothness,
    \item stability,
    \item reversibility,
    \item monotonicity,
    \item altitude/domain semantics,
    \item invariant correctness.
\end{itemize}

This completes the formal development of the runtime semantics.
